\documentclass{article}
\usepackage{amsmath, amssymb}
\usepackage{cancel}
\usepackage{tikz}
\usepackage{mdframed}
\usepackage{hyperref}
\usetikzlibrary{arrows.meta, decorations.markings}

\begin{document}

\section{Lecture 11: Plan for Class}
Last time:
\begin{enumerate}
    \item Co-dimension 1 bifurcations of flows
    \begin{enumerate}
        \item $\dot{x}=\mu-x^2$ saddle-node
        \item $\dot{x}=\mu x-x^2$ transcritical
        \item $\dot{x}=\mu x\pm x^3$ pitchfork (sub, super)-critical
    \end{enumerate}
    \item Theorem
    \begin{enumerate}
        \item eigenvalue on th imaginary axis $\left(\frac{\partial f}{\partial x}(x_0)=0\right)$
        \item nonlinear terms must not be degenerate $\left( \frac{\partial' f}{\partial x^2}(x_0)\neq 0 \right)$
        \item eigenvalues cross with nonzero speed $\left(\frac{\partial f}{\partial \mu}(x_0)\neq 0\right)$
    \end{enumerate}
\end{enumerate}

Today:
\begin{enumerate}
    \item Checkout \href{https://vanderbei.princeton.edu/nBody_animations/index.html}{this!}
    \item Using bifurcation theorems
    \item Hopf bifurcations
    \item Co-dimension 1 bifurcations of maps (maybe)
\end{enumerate}

\section{Reviewing bifurcation theorems}

Just to remind ourselves, the conditions are 
\begin{enumerate}
    \item eigenvalue on th imaginary axis $\left(\frac{\partial f}{\partial x}(x_0)=0\right)$
    \item nonlinear terms must not be degenerate $\left( \frac{\partial' f}{\partial x^2}(x_0)\neq 0 \right)$
    \item eigenvalues cross with nonzero speed $\left(\frac{\partial f}{\partial \mu}(x_0)\neq 0\right)$
\end{enumerate}

\subsection{Why condition (2)?}

Why do we need condition (2)? Let's consider the system
\[
\dot{x}=\mu x
\]
at $\mu=0$. We see that the fixed points are at $x_0=0$, unless $\mu=0$ when we get that all values of $x$ are fixed points. So the phenomena exists briefly, then disappears; hence, we don't want to consider this a bifurcation.

We can consider a less trivial example of 
\begin{align*}
    \dot{x}&= y\\
    \dot{y}&= -x + \mu y
\end{align*}
Notice that at $\mu=0$ we have 
\begin{align*}
    \dot{x}&= y\\
    \dot{y}&= -x
\end{align*}
which is just a rotation. But this rotation only exists briefly, then disappears; hence, we don't want to consider this a bifurcation.

\subsection{Why condition (3)?}

Let's consider the system 
\[
\dot{x}=\mu^2 + x^2
\]
where the fixed point at $(0,0)$ only exists for $\mu=0$. This is not a bifurcation, and we can show that this phenomena happens because $\left(\frac{\partial f}{\partial \mu}(x_0) = 0\right)$.

Perturbing our system gives us 
\[
\dot{x}=v + \mu^2 + x^2
\]
where when $v>0$, we get no bifurcations or equilibrium points. When $v<0$, we get at ``bubble'' of fixed points at $x=\pm\sqrt{-\mu^2 -v}$. I forgot what this perturbation shows us.

\section{Higher dimensions!}

\subsection{Higher dimensional saddle-node bifurcation}

\underline{Theorem 3.4.1} For 
\[
\dot{x}=f(x,\mu) \qquad x\in\mathbb{R}^n \qquad \mu\in\mathbb{R}
\]
suppose $p$ is a fixed point at $\mu_0$ such that $f(p,\mu_0)=0$. If 
\begin{enumerate}
    \item $D_x f(p,\mu_0)$ has a simple zero eigenvalue with right eigenvector $v$ (and left eigenvector $w$) and no other eigenvalues with zero real part. 
    \item $w\cdot D^2_x f(p,\mu_0)(v,v)\neq 0\implies \sum^n_{i=1}w_i\sum_{j,k=1}^{n}\frac{\partial^2 f_i}{\partial x_j \partial x_k}(p,\mu_0)v_j v_k\neq 0$
    \item $w\cdot\frac{\partial f}{\partial \mu}(p,\mu_0)\neq0\implies \sum_{i=1}^n w_i\frac{\partial f_i}{\partial \mu}(p,\mu_0)\neq 0$
\end{enumerate}
Then there is a saddle-node bifurcation at $\mu_0$ and there exists a curve of equilibrium points passing through $(p,\mu_0)$, tangent to $\mathbb{R}^n\times \left\{\mu_0\right\}$.

\subsection{Example}
Let our system be 
\begin{align*}
    \dot{x}&= y\\
    \dot{y} &= \mu - x^2 - y
\end{align*}
where we will denote our system as $f(x,y)$. At $\mu=0$ there is a fixed point at $(0,0)$. Our Jacobian is 
\[
D_x f(x,y)\Big|_{\mu=0}=\begin{bmatrix}
    0 & 1 \\ -2x & -1
\end{bmatrix}
\]
and applying it to our fixed point at $(0,0)$, we have 
\[
D_x f(0,0)\Big|_{\mu=0}=\begin{bmatrix}
    0 & 1 \\ 0 & -1
\end{bmatrix}
\]
where the eigenvalues are at $0,-1$. The eigenvectors are 
\begin{itemize}
    \item $\lambda=0$ has a right eigenvector of $v=\begin{pmatrix}
        1\\0
    \end{pmatrix}$ and a left eigenvector of $w=\begin{pmatrix}
        1 & 1
    \end{pmatrix}$
    \item $\lambda=-1$ has a right eigenvector of $v=\begin{pmatrix}
        1\\-1
    \end{pmatrix}$ and a left eigenvector of $w=\begin{pmatrix}
        0 & 1
    \end{pmatrix}$
\end{itemize}

\subsubsection{Method 1 for finding bifurcations}

To find the bifurcations, we can do 
\begin{enumerate}
    \item Find center manifold
    \item look at the 1-dimension dynamics in the center manifold
\end{enumerate}
To begin, let's diagonalize by writing 
\[
T = \begin{bmatrix}
    1 & 1 \\ 0 & -1
\end{bmatrix}\qquad T^{-1} = \begin{bmatrix}
    1 & 1 \\ 0 & -1
\end{bmatrix}
\]
where we have 
\[
\begin{pmatrix}
    x \\ y
\end{pmatrix} = T \begin{pmatrix}
    u \\ v
\end{pmatrix}
\]
Hence it follows that 
\[
\begin{pmatrix}
    \dot{u} \\ \dot{v}
\end{pmatrix} = T^{-1}\begin{pmatrix}
    0 \\ u
\end{pmatrix} + T^{-1}AT \begin{pmatrix}
    u \\ v
\end{pmatrix} + T^{-1}f\left( T\begin{pmatrix}
    u \\ v
\end{pmatrix} \right)
\]
where 
\begin{itemize}
    \item $\dot{u}=\mu -(u+v )^2$
    \item $\dot{v} = -\mu - v +(u+v )^2$
    \item $\dot{\mu}=0$
\end{itemize}
Where we can go about finding the center manifold in this diagonalized space!

Now, at $\mu=0$, we have 
\begin{align*}
    \dot{u}&= -(u+v )^2\\
    \dot{v} &=  v +(u+v )^2
\end{align*} 
where it follows that 
\[
Dg(u,v)=\begin{bmatrix}
    -2(u+v) & -2(u+v)\\2(u+v) & 1 + 2(u+v)
\end{bmatrix},\qquad Dg(0,0)=\begin{bmatrix}
    0 & 0 \\ 0 & 1
\end{bmatrix}
\]
justifying that the center manifold can be approximated as 
\[
v = h(u,\mu) = au^2 + b\mu
\]
But why do we have this $b\mu$ term? This is because our system is not actually diagonal when considering $\mu$! We have 
\[
\begin{bmatrix}
    0 & 0 & 1\\
    0 & -1 & -1 \\
    0 & 0 & 0
\end{bmatrix}
\]
where the last row and last column are the $\mu$ coefficients. 
\begin{mdframed}
\textbf{Addendum}: When performing center manifold reduction, we must be careful about 
the role of the parameter $\mu$. By treating $\mu$ as a dynamic variable 
with $\dot{\mu} = 0$, our full state vector becomes $(u, v, \mu)$, and 
the linearization of the extended system is
\[
\begin{pmatrix}\dot{u} \\ \dot{v} \\ \dot{\mu}\end{pmatrix} = 
\begin{bmatrix} 0 & 0 & 1\\ 0 & -1 & -1 \\ 0 & 0 & 0 \end{bmatrix}
\begin{pmatrix}u \\ v \\ \mu\end{pmatrix}
\]
where the linear terms are read off from the diagonalized system:
\begin{itemize}
    \item $\dot{u} = \mu - (u+v)^2$ contributes linear part $\mu$, giving row $(0,0,1)$
    \item $\dot{v} = -v - \mu + (u+v)^2$ contributes linear part $-v - \mu$, giving row $(0,-1,-1)$
    \item $\dot{\mu} = 0$ gives row $(0,0,0)$
\end{itemize}
This extended matrix has eigenvalues $0, -1, 0$, so the center subspace 
is \textbf{two-dimensional}, spanned by the $u$ and $\mu$ directions. 
Consequently, the center manifold is a surface over $(u,\mu)$, and we 
must approximate it as $v = h(u, \mu)$.

The key point is the off-diagonal $-1$ entry in position $(2,3)$, which 
reflects the fact that $\mu$ couples linearly into the $\dot{v}$ equation. 
This forces us to include a linear $\mu$ correction in our approximation, giving
\[
v = h(u,\mu) = au^2 + b\mu + O(3)
\]
Had the system been fully diagonal in $(u, v, \mu)$, this coupling would 
be absent and the $b\mu$ term could be dropped.
\end{mdframed}
Now here comes the center manifold approximation algebra. We can write
\[
\dot{v}= 2au\dot{u} + b \dot{\mu}= - \mu - (au^2 + b\mu) + \left( u+ au^2 + b\mu \right)^2
\]
where remember that $\dot{\mu}=0$. The expression above implies 
\[
2au\left(\mu - \left( u+ au^2 + b\mu \right)^2 \right) = - \mu - (au^2 + b\mu) + \left( u+ au^2 + b\mu \right)^2
\]
where the coefficients can be determined by 
\begin{enumerate}
    \item $u^2\longrightarrow 0=-a+1\implies a=1$
    \item $\mu\longrightarrow 0 = -1 - b \implies b=-1$
\end{enumerate}
Hence, our center manifold is 
\[
v = u^2 - \mu + O(3)
\]
and the dynamics in the center manifold are 
\begin{align*}
    \dot{u} &= \mu - (u+v)^2 = \mu - (u + u^2 - \mu)^2\\
    &= \mu - u^2 + O(3)
\end{align*}
which is just a saddle-node bifurcation!

\subsubsection{Method 2}

For method 2, let's try and satisfy all of the conditions of \underline{Theorem 3.4.1} to conclude that the system has a saddle-node bifurcation.

\begin{enumerate}
    \item Yes! The system has a simple zero eigenvalue with right eigenvector $v$ (and left eigenvector $w$) and no other eigenvalues with zero real part because our eigenvalues are $0,-1$ with the appropriate left and right eigenvectors.
    \item Let's check to see that the condition \[\sum^n_{i=1}w_i\sum_{j,k=1}^{n}\frac{\partial^2 f_i}{\partial x_j \partial x_k}(p,\mu_0)v_j v_k\neq 0\] holds. This is a bit of a messy calculation, but it's better than the alternative.
\end{enumerate}
We start by writing computing $D_x^2 f_1(0,0)$ and writing
\[
D_x^2 f_1(0,0) =
\begin{pmatrix}
    \dfrac{\partial^2 f_1}{\partial x^2} & \dfrac{\partial^2 f_1}{\partial x \partial y} \\[10pt]
    \dfrac{\partial^2 f_1}{\partial y \partial x} & \dfrac{\partial^2 f_1}{\partial y^2}
\end{pmatrix}
=
\begin{pmatrix}
    0 & 0 \\
    0 & 0
\end{pmatrix}
\]
where $\dot{x}=f_(x,y)=y$. This implies that 
\[
D_x^2 f_1(p, \mu)(v, v) = \sum_{j,k} \left(\frac{\partial^2 f_1}{\partial x_j \, \partial x_k}\right)_0 v_j \, v_k = 0
\]
Moving on to computing $D_x^2 f_2(0,0)$, we can write
\[
D_x^2 f_2(0,0) =
\begin{pmatrix}
    -2 & 0 \\
    0  & 0
\end{pmatrix}
\]
where $\dot{y} =f_2(x,y)= \mu - x^2 - y$. This implies that 
\[
D_x^2 f_2(p, \mu)(v, v) = v^T
\begin{bmatrix}
-2 & 0 \\
0  & 0
\end{bmatrix}
v
= (1 \quad 0)
\begin{bmatrix}
-2 & 0 \\
0  & 0
\end{bmatrix}
\begin{pmatrix} 1 \\ 0 \end{pmatrix}
= -2
\]
Now we can put it all together and write
\[
D^2_x f(p,\mu_0)(v,v) = \begin{pmatrix}
    D_x^2 f_1(p, \mu)(v, v) \\ D^2_x f(p,\mu_0)(v,v)
\end{pmatrix} = \begin{pmatrix}
    0 \\ -2
\end{pmatrix}
\]
which implies that
\[
w\cdot D^2_x f(p,\mu_0)(v,v)=w \cdot\begin{pmatrix} 0 \\ -2 \end{pmatrix}= (1 \quad 1)\begin{pmatrix} 0 \\ -2 \end{pmatrix}= -2 \neq 0
\]

\begin{enumerate}
    \item[2.] Now we can confirm that \[\sum^n_{i=1}w_i\sum_{j,k=1}^{n}\frac{\partial^2 f_i}{\partial x_j \partial x_k}(p,\mu_0)v_j v_k\neq 0\]
    \item[3.] Let's check to see that the condition \[w\cdot\frac{\partial f}{\partial \mu}(p,\mu_0)\neq0\] holds! This calculation is much nicer.
\end{enumerate}
It should be fairly clear to see that
\[
\frac{\partial f}{\partial \mu} = \begin{pmatrix}
    0\\1
\end{pmatrix}\]
which implies that
\[w\cdot \frac{\partial f}{\partial \mu}= \begin{pmatrix}
    1 & 1
\end{pmatrix}\begin{pmatrix}
    0\\1
\end{pmatrix}=1\neq 0
\]
\begin{enumerate}
    \item[3.] Now we can confirm that \[w\cdot\frac{\partial f}{\partial \mu}(p,\mu_0)\neq0\]
\end{enumerate}
Therefore, we have a saddle node bifurcation!

\subsubsection{Note}
If we modify the equations to 
\begin{align*}
    \dot{x}&= y\\
    \dot{y}&= \mu x -x^2 - y
\end{align*}
then we will get a transcritical bifurcation.

\section{Hopf bifurcation}

Suppose we have the system
\begin{align*}
    \dot{x}&= \mu x - y + ax\left(x^2 + y^2\right)^2 - by\left(x^2 + y^2\right)^2 \\
    dot{y}&= x + \mu y + ay\left(x^2 + y^2\right)^2 - bx\left(x^2 + y^2\right)^2 \\
\end{align*}
where we have an equilibrium point at $(x,y)=(0,0)$. Performing a linearization, we have 
\[
Df(0,0)=\begin{bmatrix}
    \mu & -1 \\ 1 & \mu
\end{bmatrix}
\]
where the eigenvalues are $\mu \pm i$. 

At $\mu=0$, we have a Hopf bifurcation. Looking at the linearization, we have 
\[
Df(0,0)=\begin{bmatrix}
    0 & -1 \\ 1 & 0
\end{bmatrix}
\]
and our system becomes (as we saw from the normal form)
\begin{align*}
    \dot{x}&= - y + ax\left(x^2 + y^2\right)^2 - by\left(x^2 + y^2\right)^2 \\
    dot{y}&= x  + ay\left(x^2 + y^2\right)^2 - bx\left(x^2 + y^2\right)^2 \\
\end{align*}

\subsection{Polar coordinate time}

As a reminder, polar coordinate transformations look like
\begin{align*}
    x &= r\cos\theta \implies \dot{x}=\dot{r}\cos\theta - r\sin\theta\,\dot{\theta} \text{ denoted as (A)} \\
    y &= r\sin\theta \implies \dot{y}=\dot{r}\sin\theta + r\cos\theta\,\dot{\theta} \text{ denoted as (B)}
\end{align*}
Embarking on our adventure to eliminate terms, we can start by noting that
\begin{align*}
    (A)\cos\theta + (B)\sin\theta &= \left(\dot{r}\cos\theta - r\sin\theta\,\dot{\theta}\right)\cos\theta + \left(\dot{r}\sin\theta + r\cos\theta\,\dot{\theta}\right)\sin\theta\\
    & = \dot{r}\cos^2\theta - \cancel{r\sin\theta\,\cos\theta\,\dot{\theta}} + \dot{r}\sin^2\theta + \cancel{r\sin\theta\,\cos\theta\,\dot{\theta}}\\
    &= \dot{r}\cos^2\theta+ \dot{r}\sin^2\theta\\
    &=\dot{r}\left( \cos^2\theta+ \sin^2\theta \right)\\
    &=\dot{r}
\end{align*}
Hence we are justified in writing
\begin{align*}
\dot{r} &= \dot{x}\cos\theta + \dot{y}\sin\theta \\
        &= \left(-y + axr^2 - \cancel{byr^2}\right)\cos\theta + \left(x + ayr^2 + \cancel{bxr^2}\right)\sin\theta \\
        &= -r\sin\theta\cos\theta + ax^2r + r\cancel{\cos\theta\sin\theta} + a y^2 r \\
        &= ar(x^2 + y^2) \\
        &= ar^3
\end{align*}
Hence we have
\[
\boxed{\dot{r} = ar^3}
\]
where it follows that
\[
a < 0: \text{ stable} \qquad a > 0: \text{ unstable} \qquad a = 0: r = \text{const.}
\]
Now we could compute $\dot{\theta}$, but that's not too relevant. Let's just move on.

\subsection{A theorem!}

\underline{Theorem} If $Df(x_0, \mu_0)$ has a simple pair of complex eigenvalues crossing the imaginary axis, then there exists a smooth curve of equilibria $x(\mu)$ with $x(\mu_0)=x_0$ and a 2d surface of periodic orbits tangent to the center eigenspace. If 
\begin{itemize}
    \item $a<0$ then we have a supercritical Hopf bifurcation with a stable limit cycle,
    \item $a>0$ then we have a subcritical Hopf bifurcation with an unstable limit cycle.
\end{itemize}
There's a cool image here of the Hopf bifurcation that I should include.

\subsection{A formula for coefficients in the normal form }

As seen in page 152, the formula for $a$ that appears in our normal form is 
\begin{align*}
    a &= \frac{1}{16}\left[ f_{xxx} + f_{xyy} + g_{xxy} + g_{yyy} \right]\\
    &+ \frac{1}{16\omega}\left[ f_{xy}\left(f_{xx}+f_{yy}\right) - g_{xy}\left( g_{xx} + g_{yy} \right) -f_{xx}g_{xx} + f_{yy}g_{yy} \right]
\end{align*}
where 
\begin{align*}
    \dot{x}&= f(x,u,\mu)\\
    \dot{y}&= g(x,y,\mu)
\end{align*}
and
\[
Df(x_0,\mu)=\begin{bmatrix}
    \mu & -\omega\\ \omega & \mu
\end{bmatrix}
\]

\subsubsection{An example}
Let's say that we have the system
\begin{align*}
    \dot{x}&= -\omega y + \alpha y^2 - \beta x^2 y = f\\
    \dot{y}&= \omega x - \gamma y^2 + \delta xy-y^3=g
\end{align*}
Before we go about approximating $a$, let's first note the following terms 
\begin{gather*}
    f_{xx} =0\\
    f_{xy}=0\\
    f_{yy}=2\alpha \\
    f_{xxx}=f_{xyy} = 0\\
    g_xx=0\\
    g_{xy}=\delta\\
    g_{yy}=-2\gamma\\
    g_{xxy}=0\\
    g_{yyy}=-6
\end{gather*}
Maybe there is a mistake, but who cares. This allows us to write 
\[
a=\frac{1}{8}\left(\frac{\gamma}{\omega}\left(\delta - 2\alpha\right)-3\right)
\]
Pretty cool, right?

\end{document}